\clearpage\begingroup\SourceLeaf{437b}\ArithmeticHead{106}{L}{STÄCKEL, GAUSS ALS GEOMETER.}
\renewcommand{\thefootnote}{\arabic{footnote})}\setcounter{footnote}{0}
\begin{center}
\textbf{32.}\par\smallskip
\textls[100]{Weitere Untersuchungen über krumme Flächen.}
\end{center}
\noindent\hspace*{\parindent}In der Selbstanzeige der \textit{Disq. gen.} sagt \textsc{Gauss}, der Zweck der Abhandlung
sei, neue Gesichtspunkte für die Lehre von den krummen Flächen zu
eröffnen und einen Teil der neuen Wahrheiten, die dadurch zugänglich werden,
\clearpage\SourceLeaf{438}\ArithmeticHead{107}{R}{DIE ALLGEMEINE LEHRE VON DEN KRUMMEN FLÄCHEN.}\setcounter{footnote}{0}
\noindent zu entwickeln (W. IV, S. 341). Dass dort nur ein Teil der Ergebnisse, zu
denen er gelangt war, dargestellt ist, wird auch in den Briefen an \textsc{Bessel}, \textsc{Olbers}
und \textsc{Schumacher} ausgesprochen, die in der vorhergehenden Nummer angeführt
sind, ja es wird einmal geradezu eine zweite Abhandlung über die krummen
Flächen in Aussicht gestellt (Brief an \textsc{Olbers} vom 1. März 1827, W. IX, S. 377).
\par
Auch in den \textit{Disq. gen.} finden sich Andeutungen über weitergehende
Untersuchungen. So werden im art. 6 Erörterungen über die allgemeinste
Auffassung des Inhalts von Figuren auf eine andere Gelegenheit verschoben.
Ferner unterscheidet \textsc{Gauss} im art. 13 zwischen den Eigenschaften einer
krummen Fläche, die von ihrer gerade angenommenen Form abhängen, und
jenen, die erhalten bleiben, in welche Form die Fläche auch gebogen wird.
Hierfür nennt er das Krümmungsmass, die Lehre von den kürzesten Linien
und einiges andere, dessen Behandlung er sich vorbehalte.
\par
Zu den Gegenständen, die im art. 13 gemeint sind, gehört vor allem die
»Seitenkrümmung« von Kurven auf krummen Flächen, die \textsc{Gauss} schon in der
Zeit zwischen 1822 und 1825 eingehend untersucht hatte (W. VIII, S. 386—
395). Eine solche Kurve besitzt zunächst eine \textls[100]{absolute Krümmung}, die
durch den reziproken Wert des auf die übliche Art erklärten Krümmungshalbmessers
gegeben wird. Wenn man aber den Krümmungshalbmesser in
zwei Komponenten, nach der Flächennormale und senkrecht dazu, zerlegt, so
werden in deren reziproken Werten die Masse der \textls[100]{Normalkrümmung} und
der \textls[100]{Seitenkrümmung} gewonnen. Die kürzesten Linien auf der Fläche
haben die Eigenschaft, dass ihr Krümmungshalbmesser in die zugehörige
Flächennormale fällt, und ihnen kommt daher die Seitenkrümmung Null zu,
sie entsprechen auch in dieser Hinsicht den geraden Linien der Ebene, und
in einer Geometrie der auf einer krummen Fläche liegenden Figuren, bei der
an die Stelle der Geraden die Kürzesten treten, ist bei einer Kurve die \textls[100]{relative Krümmung},
das heisst das Verhältnis des geodätischen Kontingenzwinkels
zum Linienelement der Kurve, gleich der Seitenkrümmung zu setzen.
Bald nach dem Erscheinen der \textit{Disq. gen.} hat übrigens \textsc{Minding} ähnliche Auffassungen
veröffentlicht\footnote{\textsc{F. Minding}, \textit{Über die Kurven kürzesten Perimeters auf krummen Flächen}, \textsc{Crelles} Journal, Bd. 5, 1830, S. 297.}.
\par
Im Laufe der Untersuchung überträgt \textsc{Gauss} den Namen der Seitenkrümmung
\par\smallskip\noindent\hfill{\small 14\textsuperscript{*}}\par
\clearpage\SourceLeaf{439}\ArithmeticHead{108}{L}{STÄCKEL, GAUSS ALS GEOMETER.}\setcounter{footnote}{0}
\noindent auf das über die Kurve erstreckte Integral der ursprünglichen
Seitenkrümmung. Er hatte dabei wohl die Verallgemeinerung des Satzes von
der Winkelsumme des geodätischen Dreiecks im Auge, die später von \textsc{Bonnet}
angegeben worden ist\footnote{\textsc{O. Bonnet}, \textit{Mémoires sur la théorie générale des surfaces}, Journal de l’école polytechnique, t. 19, cah. 32, 1848, S. 131.}. Hiernach ist die Gesamtkrümmung eines beliebigen
auf einer krummen Fläche liegenden Dreiecks gleich dem Unterschiede der
Winkelsumme gegen zwei Rechte, vermindert um das über die Begrenzung
erstreckte Integral der Seitenkrümmung (im ursprünglichen Sinne des Wortes).
\par
Die Erklärung der kürzesten Linien als der Kurven von der Seitenkrümmung
Null ist auch insofern wichtig, als die Rechnungen, die \textsc{Gauss}
daran anschliesst, einen Einblick in die Kunstgriffe gewähren, die ihn zu den
eleganten Formeln im art. 22 der \textit{Disq. gen.} geführt haben.
\par
Ob \textsc{Gauss} die Geometrie der Figuren auf einer krummen Fläche noch
weiter ausgebaut, ob er im besonderen den Zusammenhang zwischen der Geometrie
auf den Flächen konstanten Krümmungsmasses und der nichteuklidischen
Geometrie der Ebene erkannt hat, ist nicht mit Sicherheit zu entscheiden.
Nahe genug musste er für jemand liegen, der schon 1794 wusste, dass dort
das Verhältnis des Dreiecksinhaltes zu der Abweichung der Winkelsumme von
zwei Rechten eine Konstante ist (W. VIII, S. 266). Auch die Bemerkungen,
dass die Untersuchungen über die krummen Flächen so vielfach in andere Teile
der Mathematik eingriffen (Brief an \textsc{Bessel} vom 20. November 1826, W. IX,
S. 362), dass sie tief in die Metaphysik der Raumlehre eingreifen (Brief an \textsc{Hansen}
vom 11. Dezember 1825, Gauss-Archiv) in Verbindung mit der Tatsache, dass
\textsc{Gauss} bald nach Vollendung der \textit{Disq. gen.} die Untersuchungen über die
Grundlagen der Geometrie wieder aufgenommen hat (Brief an \textsc{Bessel} vom
27. Januar 1829, W. VIII, S. 200), lassen sich zu Gunsten einer solchen Annahme
geltend machen. Ferner wird in einer Aufzeichnung aus dem Jahre
1846 (W. VIII, S. 257) die einer nichteuklidischen Geometrie eigentümliche
absolute Konstante mit \(k\) bezeichnet, wo \(k\) die Quadratwurzel aus dem Krümmungsmass
bedeuten würde. Bemerkenswert ist auch eine Wendung in einem
aus demselben Jahre 1846 stammenden Briefe an \textsc{Gerling}: »Der Satz, den
Ihnen Herr \textsc{Schweikart} erwähnt hat, dass in \textls[100]{jeder} Geometrie die Summe
aller äussern Polygonwinkel von \(360^0\) um eine Grösse verschieden ist (nämlich
\clearpage\SourceLeaf{440}\ArithmeticHead{109}{R}{DIE ALLGEMEINE LEHRE VON DEN KRUMMEN FLÄCHEN.}\setcounter{footnote}{0}
\noindent grösser als \(360^0\) in der Astralgeometrie, wie \textsc{Schweikart} sie aufgefasst hat),
welche dem Flächeninhalt proportional ist, ist der erste, gleichsam an der
Schwelle liegende Satz der Theorie, den ich schon im Jahr 1794 als notwendig
erkannte« (Brief vom 2. Oktober 1846, W. VIII, S. 266); \textsc{Gauss} unterscheidet
also die Auffassung \textsc{Schweikarts} von der seinigen, bei der in \textls[100]{jedem}
Falle die Winkelsumme des Dreiecks von \(180^0\) verschieden ist, sodass bei ihm
neben die Geometrie, bei der die Winkelsumme kleiner als \(180^0\) ist, noch
eine zweite tritt, bei der die Winkelsumme grösser als \(180^0\) wird. Wenn
man beachtet, wie vorsichtig \textsc{Gauss} bei solchen Andeutungen zu Werk ging
(vgl. S. 9), so wird man auch auf diese Stelle Gewicht zu legen haben.
\par
Schliesslich verdient erwähnt zu werden, dass in einer spätestens 1827
niedergeschriebenen Aufzeichnung die durch Drehung der Traktrix entstehende
krumme Fläche negativen konstanten Krümmungsmasses (Pseudosphäre) als
das »Gegenstück der Kugel« bezeichnet wird (W. VIII, S. 265). \textsc{Gauss} erwähnt
die Pseudosphäre im Zusammenhang mit der Verbiegung von Drehflächen
in Drehflächen. Aber noch mehr, die von ihm aufgestellten Formeln
führen zu dem Satz, dass bei der Pseudosphäre (und nur bei ihr) alle diese
Drehflächen einander kongruent sind, und hierin liegt, dass man ein geodätisches
Dreieck, unter Bewahrung dieser Eigenschaft, auf der Pseudosphäre
ebenso verschieben kann wie ein sphärisches Dreieck auf der Kugel. Hat
\textsc{Gauss} deshalb den Namen »Gegenstück der Kugel« gewählt? Jedenfalls hat
er den krummen Flächen von negativem konstanten Krümmungsmass seine
Aufmerksamkeit zugewendet. In den schönen Untersuchungen, die \textsc{Minding},
angeregt durch die \textit{Disq. gen.}, angestellt hat, sind auch diese Ergebnisse über
die Biegung der Drehflächen und über die Pseudosphäre enthalten\footnote{\textsc{F. Minding}, \textit{Über die Biegung gewisser Flächen}, \textsc{Crelles} Journal, Bd. 18, 1838, S. 367; \textit{Wie sich entscheiden lässt, ob zwei gegebene krumme Flächen auf einander abwickelbar seien oder nicht}, ebenda, Bd. 19, 1839, S. 378; \textit{Über die kürzesten Linien krummer Flächen}, ebenda, Bd. 20, 1840, S. 324.}.
\par
Auf die Biegung krummer Flächen bezieht sich auch eine wahrscheinlich
Ende 1826 niedergeschriebene kurze Bemerkung, in der \textsc{Gauss} die Beziehung,
die bei zwei Biegungsflächen zwischen den sphärischen Abbildungen mittels
paralleler Normalen besteht, zu einem Ansatz für die Lösung des allgemeinen
Problems der Abwicklung krummer Flächen auf einander benutzt, der erst
im Jahre 1900 aus dem Nachlass im achten Bande der Werke (S. 447—448)
\clearpage\SourceLeaf{441a}\ArithmeticHead{110}{L}{STÄCKEL, GAUSS ALS GEOMETER.}\setcounter{footnote}{0}
\noindent veröffentlicht worden ist. Es wäre zu wünschen, dass dieser Gedanke, der
\textsc{Gauss} eigentümlich ist, vollständig durchgeführt würde.
\par
Zum Schluss sei noch berichtet, dass die philosophische Fakultät der
Universität Göttingen im Jahre 1830 auf Veranlassung von \textsc{Gauss} die Preisfrage
stellte: Determinetur inter lineas duo puncta jungentes ea, quae circa
datum axem revoluta gignat superficiem minimam. Sie wurde von seinem
Landsmann, Schüler und späteren Mitarbeiter auf der Sternwarte, \textsc{Goldschmidt},
beantwortet, dem auch der Preis zugefallen ist\footnote{\textsc{B. Goldschmidt}, \textit{Determinatio superficiei minimae rotatione curvae data duo puncta jungentis circa datum axem ortae}, Göttingen 1831.}.
\clearpage\endgroup
